\chapter{Ricerca e ordinamento su array in Java}
\label{chap:ricerca-ordinamento-java}

Gli algoritmi introdotti nel \cref{chap:ricerca-ordinamento} vengono
realizzati su array di \texttt{int} mediante i metodi statici della
classe \texttt{ArrayAlgs} (\cref{sec:fi-java-array-algoritmi}).
Si elaborano tutti gli elementi dell'array, senza un parametro
separato per la lunghezza. Le versioni per oggetti confrontabili
sono trattate nella \cref{sec:fi-java-interfacce-arrayalgs}.
Le analisi assumono il modello di costo
unitario della \cref{sec:fi-complessita-modello}.\footnote{Materiale
di riferimento: S. Canazza, \emph{Java: Ricerca e Ordinamento},
\nolinkurl{2026_FI_13_canazza_Java__RicercaOrdinamento.key.pdf}.}

\section{Ricerca lineare}
\label{sec:fi-java-alg-lineare}

\texttt{linearSearch()} visita l'array dall'indice zero e restituisce
la posizione della prima occorrenza del valore cercato, oppure
\(-1\) se il valore è assente. Non richiede che l'array sia ordinato
(\cref{sec:fi-ricerca-lineare}):

\noindent
\begin{minipage}{\linewidth}
\begin{unipdcode}[language=Java]
public static int linearSearch(int[] v, int value) {
    if (v == null) throw new IllegalArgumentException();
    for (int i = 0; i < v.length; i++)
        if (v[i] == value) return i;
    return -1;
}
\end{unipdcode}

\end{minipage}

Prima dell'iterazione di indice \(i\), nessuna posizione precedente
contiene il valore richiesto; il primo confronto riuscito individua
quindi la prima occorrenza. Per \(n\ge1\), il caso migliore richiede
\(\Theta(1)\) tempo e il peggiore \(\Theta(n)\); il costo medio è
\(\Theta(n)\) assumendo una prima occorrenza uniformemente
distribuita fra le posizioni. Lo spazio ausiliario è
\(\Theta(1)\). Con un array vuoto il metodo restituisce \(-1\).

\section{Ricerca binaria}
\label{sec:fi-java-alg-binaria}

La ricerca binaria richiede un array ordinato in senso
\emph{non decrescente}. A ogni passo confronta il valore cercato
con l'elemento centrale dell'intervallo candidato ed esclude la
metà incompatibile con l'ordinamento
(\cref{sec:fi-ricerca-binaria}).

\subsection{Versione iterativa}
\label{sec:fi-java-alg-binaria-iterativa}

\texttt{binarySearch()} mantiene due indici, \texttt{left} e
\texttt{right}, che delimitano l'intervallo ancora da esaminare:

\noindent
\begin{minipage}{\linewidth}
\begin{unipdcode}[language=Java]
public static int binarySearch(int[] v, int value) {
    if (v == null) throw new IllegalArgumentException();
    int left = 0, right = v.length - 1;
    while (left <= right) {
        int mid = left + (right - left) / 2;
        if (v[mid] == value) return mid;
        if (v[mid] < value) left = mid + 1;
        else right = mid - 1;
    }
    return -1;
}
\end{unipdcode}

\end{minipage}

Se \texttt{v[mid]} è minore del valore cercato, si aggiorna
\texttt{left}; in caso contrario si riduce \texttt{right}.
L'intervallo si restringe a ogni iterazione e diventa vuoto
quando \lstinline|left > right|, condizione che determina il ritorno
di \(-1\).

\subsection{Versione ricorsiva}
\label{sec:fi-java-alg-binaria-ricorsiva}

La versione ricorsiva passa gli estremi dell'intervallo inclusivo
al metodo ausiliario \texttt{binSearch()}:

\noindent
\begin{minipage}{\linewidth}
\begin{unipdcode}[language=Java]
public static int binarySearchRecursive(int[] v, int value) {
    if (v == null) throw new IllegalArgumentException();
    return binSearch(v, 0, v.length - 1, value);
}

private static int binSearch(int[] v, int from, int to, int value) {
    if (from > to) return -1;
    int mid = from + (to - from) / 2;
    if (v[mid] == value) return mid;
    if (v[mid] < value)
        return binSearch(v, mid + 1, to, value);
    return binSearch(v, from, mid - 1, value);
}
\end{unipdcode}

\end{minipage}

Ogni chiamata ricorsiva conserva soltanto una delle due metà.
Il caso base corrisponde a un intervallo vuoto o a un elemento
uguale al valore cercato. Entrambe le versioni restituiscono
un indice valido oppure \(-1\); in presenza di duplicati non
è garantita la prima occorrenza.

Per \(n\ge2\), il tempo è \(\Theta(1)\) nel caso migliore e
\(\Theta(\log n)\) nel caso peggiore. Con valori distinti e ricerca
riuscita uniformemente distribuita, anche il caso medio è
\(\Theta(\log n)\). La versione iterativa occupa
\(\Theta(1)\) spazio ausiliario, quella ricorsiva
\(O(\log n)\) per la pila delle chiamate. L'ordinamento iniziale
rimane una precondizione per entrambe.

\section{Ordinamento per selezione}
\label{sec:fi-java-alg-selezione}

\texttt{selectionSort()} cerca la posizione del minimo nel tratto
ancora da ordinare e lo scambia con il primo elemento di tale
tratto. Si usano i metodi ausiliari \texttt{findMinPos()} e
\texttt{swap()}:

\noindent
\begin{minipage}{\linewidth}
\begin{unipdcode}[language=Java]
public static void selectionSort(int[] v) {
    if (v == null) throw new IllegalArgumentException();
    for (int i = 0; i < v.length - 1; i++) {
        int minPos = findMinPos(v, i);
        if (minPos != i) swap(v, i, minPos);
    }
}
\end{unipdcode}

\end{minipage}

\noindent
\begin{minipage}{\linewidth}
\begin{unipdcode}[language=Java]
private static int findMinPos(int[] v, int from) {
    int pos = from;
    for (int i = from + 1; i < v.length; i++)
        if (v[i] < v[pos]) pos = i;
    return pos;
}

private static void swap(int[] v, int i, int j) {
    int temp = v[i];
    v[i] = v[j];
    v[j] = temp;
}
\end{unipdcode}

\end{minipage}

Sul suffisso di lunghezza \(n-i\), \texttt{findMinPos()} richiede
\(\Theta(n-i)\) tempo e \(\Theta(1)\) spazio; \texttt{swap()}
richiede \(\Theta(1)\) per entrambe le risorse.

Prima dell'iterazione di indice \(i\), i primi \(i\) elementi sono
ordinati e contengono i valori minori dell'array; la ricerca
del minimo nella parte rimanente mantiene questa proprietà.
Le scansioni eseguono in totale
\[
  (n-1)+(n-2)+\cdots+1=\frac{n(n-1)}{2}
\]
confronti: il tempo è \(\Theta(n^2)\) nel caso migliore, medio e
peggiore. Il metodo ordina l'array ricevuto, usa \(\Theta(1)\)
spazio ausiliario e non garantisce la stabilità
(\cref{sec:fi-ordinamento-selezione}).

\section{Ordinamento per inserimento}
\label{sec:fi-java-alg-inserimento}

\texttt{insertionSort()} mantiene ordinato un prefisso e inserisce
ciascun elemento nella posizione corretta, spostando verso destra
i valori maggiori:

\noindent
\begin{minipage}{\linewidth}
\begin{unipdcode}[language=Java]
public static void insertionSort(int[] v) {
    if (v == null) throw new IllegalArgumentException();
    for (int i = 1; i < v.length; i++) {
        int temp = v[i];
        int j = i;
        while (j > 0 && v[j - 1] > temp) {
            v[j] = v[j - 1];
            j--;
        }
        v[j] = temp;
    }
}
\end{unipdcode}

\end{minipage}

All'inizio dell'iterazione di indice \(i\), il prefisso di
lunghezza \(i\) è ordinato. Il ciclo interno sposta i valori maggiori
di \texttt{temp}; l'inserimento completa il prefisso ordinato di
lunghezza \(i+1\). Se l'array è già ordinato, il tempo è
\(\Theta(n)\). Nel caso peggiore, gli spostamenti sono
\(n(n-1)/2\), quindi il tempo è \(\Theta(n^2)\); lo stesso ordine
vale in media per permutazioni uniformi di valori distinti.
Lo spazio ausiliario è \(\Theta(1)\). Il confronto stretto
\lstinline|>| evita di spostare elementi uguali: il metodo è stabile
(\cref{sec:fi-ordinamento-inserimento}).

\section{Ordinamento per fusione}
\label{sec:fi-java-alg-fusione}

\texttt{mergeSort()} divide l'array in due parti, ordina entrambe
ricorsivamente e fonde i risultati. La suddivisione è mostrata
nella \cref{fig:fi-merge-sort-suddivisione}.

\noindent
\begin{minipage}{\linewidth}
\begin{unipdcode}[language=Java]
public static void mergeSort(int[] v) {
    if (v == null) throw new IllegalArgumentException();
    if (v.length < 2) return;
    int mid = v.length / 2;
    int[] left = new int[mid];
    int[] right = new int[v.length - mid];
    System.arraycopy(v, 0, left, 0, mid);
    System.arraycopy(v, mid, right, 0, right.length);
    mergeSort(left);
    mergeSort(right);
    merge(v, left, right);
}
\end{unipdcode}

\end{minipage}

Il metodo ausiliario \texttt{merge()} copia nell'array di
destinazione, in ordine, il minore fra i primi elementi non ancora
consumati delle due parti. Terminata una parte, copia gli
elementi residui dell'altra:

\noindent
\begin{minipage}{\linewidth}
\begin{unipdcode}[language=Java]
private static void merge(int[] v, int[] left, int[] right) {
    int i = 0, l = 0, r = 0;
    while (l < left.length && r < right.length) {
        if (left[l] <= right[r]) v[i++] = left[l++];
        else v[i++] = right[r++];
    }
    while (l < left.length) v[i++] = left[l++];
    while (r < right.length) v[i++] = right[r++];
}
\end{unipdcode}

\end{minipage}

Il metodo \texttt{merge()} trasferisce complessivamente \(n\) elementi:
richiede \(\Theta(n)\) tempo e \(\Theta(1)\) spazio ausiliario
proprio, escludendo gli array ricevuti.
La fusione conserva un prefisso ordinato contenente gli elementi
minori già elaborati. Le chiamate ricorsive ordinano le due metà;
la fusione produce quindi l'intero array ordinato. Il confronto
\lstinline|<=| sceglie per primo un elemento della metà sinistra
in caso di uguaglianza, conservando la stabilità.

La suddivisione e la fusione richiedono \(\Theta(n)\) operazioni
non ricorsive per ogni livello. La relazione
\(T(n)=2T(n/2)+\Theta(n)\), a meno degli arrotondamenti,
fornisce tempo \(\Theta(n\log n)\) in tutti i casi
(\cref{sec:fi-teorema-maestro}). Le copie temporanee occupano
\(\Theta(n)\) spazio ausiliario al picco, oltre alla pila di
profondità \(\Theta(\log n)\); il costo spaziale complessivo
resta \(\Theta(n)\).

\section{Ordinamento rapido}
\label{sec:fi-java-alg-rapido}

\texttt{quickSort()} ordina ricorsivamente i due intervalli
ottenuti mediante \texttt{partition()}. Il pivot è l'ultimo
elemento dell'intervallo considerato
(\cref{sec:fi-ordinamento-rapido}).

\noindent
\begin{minipage}{\linewidth}
\begin{unipdcode}[language=Java]
public static void quickSort(int[] v) {
    if (v == null) throw new IllegalArgumentException();
    quickSort(v, 0, v.length - 1);
}

private static void quickSort(int[] v, int left, int right) {
    if (left >= right) return;
    int pivotPos = partition(v, left, right);
    quickSort(v, left, pivotPos - 1);
    quickSort(v, pivotPos + 1, right);
}
\end{unipdcode}

\end{minipage}

Durante la partizione, l'indice \texttt{i} delimita gli elementi
non maggiori del pivot. Ogni elemento che soddisfa tale condizione
viene collocato nella parte sinistra mediante \texttt{swap()},
introdotto nella \cref{sec:fi-java-alg-selezione}. Lo scambio
finale colloca il pivot nella posizione definitiva:

\noindent
\begin{minipage}{\linewidth}
\begin{unipdcode}[language=Java]
private static int partition(int[] v, int left, int right) {
    int pivot = v[right];
    int i = left - 1;
    for (int j = left; j < right; j++) {
        if (v[j] <= pivot) {
            i++;
            swap(v, i, j);
        }
    }
    swap(v, i + 1, right);
    return i + 1;
}
\end{unipdcode}

\end{minipage}

A partizione conclusa, i valori a sinistra del pivot non lo
superano e quelli a destra sono maggiori. Le due chiamate
ricorsive operano su intervalli strettamente più piccoli e
completano l'ordinamento.

La partizione costa \(\Theta(n)\). Con divisioni bilanciate il
tempo è \(\Theta(n\log n)\); nel caso medio, per permutazioni
uniformi di elementi distinti, è ancora \(\Theta(n\log n)\).
Con partizioni ripetutamente sbilanciate, per esempio su un
array già ordinato o con tutti gli elementi uguali, scegliendo
l'ultimo elemento come pivot, il tempo diventa \(\Theta(n^2)\).

\texttt{partition()} richiede \(\Theta(1)\) spazio aggiuntivo.
La pila ricorsiva occupa \(O(\log n)\) spazio con partizioni
bilanciate e in media nelle ipotesi precedenti, fino a
\(O(n)\) nel caso peggiore. Il metodo modifica l'array
ricevuto e non è stabile.

L'adattamento degli algoritmi di ordinamento a oggetti
confrontabili è introdotto nella
\cref{sec:fi-java-ereditarieta-ordinamento-oggetti}.
